\chapter{Conclusioni}

\section{Riepilogo Punti di Forza}

Open Data Chunker si distingue per i seguenti punti di forza:

\begin{enumerate}
    \item \textbf{Robustezza}
    \begin{itemize}
        \item Sanitizzazione automatica XML malformati
        \item Fault tolerance con logging granulare
        \item Recovery automatico da errori di parsing
    \end{itemize}
    
    \item \textbf{Performance}
    \begin{itemize}
        \item Parsing streaming memory-efficient
        \item Parallelismo multi-processo scalabile
        \item Output Parquet con compressione 10x
    \end{itemize}
    
    \item \textbf{Riproducibilità}
    \begin{itemize}
        \item Containerizzazione Docker completa
        \item Ambiente consistente su qualsiasi piattaforma
        \item Integrazione CI/CD ready
    \end{itemize}
    
    \item \textbf{Usabilità}
    \begin{itemize}
        \item CLI intuitiva con Click
        \item Query SQL interattive con DuckDB
        \item Export flessibile in formati standard
    \end{itemize}
\end{enumerate}

\begin{figure}[h]
\centering
\begin{tikzpicture}[scale=0.9]
    \foreach \angle/\label/\col in {
        90/Robustezza/primaryblue,
        0/Performance/accentgreen,
        270/Riproducibilità/warnorange,
        180/Usabilità/primaryblue!60
    } {
        \draw[\col, ultra thick] (0,0) -- (\angle:2.5);
        \node at (\angle:3.2) {\textbf{\label}};
        \filldraw[\col!50] (\angle:2) circle (0.15);
    }
    \draw[gray, dashed] (0,0) circle (1.5);
    \draw[gray, dashed] (0,0) circle (2.5);
\end{tikzpicture}
\caption{I quattro pilastri di Open Data Chunker}
\end{figure}

\section{Possibili Evoluzioni}

La soluzione può essere estesa in diverse direzioni:

\begin{description}
    \item[Incremental Processing] Supporto per elaborazione differenziale di soli nuovi file, evitando re-processing completo.
    
    \item[Schema Evolution] Gestione automatica di cambiamenti nello schema XML sorgente con migrazione trasparente.
    
    \item[Data Quality] Integrazione di framework di data quality (es. Great Expectations) per validazione automatica.
    
    \item[Orchestrazione] Integrazione con Apache Airflow o Prefect per scheduling e monitoring.
    
    \item[API REST] Esposizione dei dati via API per integrazioni applicative.
\end{description}

\section{Considerazioni Finali}

Open Data Chunker dimostra come un'architettura ben progettata possa trasformare dataset pubblici complessi in risorse analitiche accessibili. La combinazione di:

\begin{itemize}
    \item Modern Python data stack (Polars, DuckDB, PyArrow)
    \item Pattern di parallelismo robusti
    \item Containerizzazione per portabilità
\end{itemize}

...rende la soluzione adatta sia per analisi esplorative occasionali che per pipeline di produzione.

\vspace{1cm}

\begin{center}
\textit{``La qualità dei dati aperti è tanto importante quanto la loro disponibilità.''}
\end{center}
